\documentclass[12pt]{article}
\usepackage[utf8]{inputenc}
\usepackage[T1]{fontenc}
\usepackage[margin=1in]{geometry}
\usepackage{setspace}
\usepackage{xcolor}
\usepackage[colorlinks=true, urlcolor=blue!60!black, linkcolor=blue!60!black]{hyperref}
\setlength{\parindent}{0em}
\setlength{\parskip}{0.45em}
\renewcommand{\baselinestretch}{1.0}
\begin{document}
\begin{center}
{\Large \textit{Ex parte McCardle}, 74 U.S. (7 Wall.) 506 (1869)\par}
\vspace{0.3em}
{\normalsize Case Brief --- One-Page Edition\par}
\end{center}

\textbf{Memory Jogger.} Congress strips a pending Reconstruction habeas appeal (\href{https://supreme.justia.com/cases/federal/us/74/506/}{74 U.S.\ 506}).

\textbf{Facts.} William McCardle, a Mississippi editor who denounced Reconstruction, was held by the military for trial by military commission. He sought habeas corpus under the 1867 Habeas Corpus Act, which allowed appeals to the Supreme Court.

\textbf{Procedural History.} After the Court heard argument on his appeal, Congress repealed the 1867 appeal provision over President Johnson's veto.

\textbf{Votes.} Dismissed, 8--0 (Chase for a unanimous Court).

\textbf{Holding.} Under the Exceptions Clause, Congress may repeal appellate jurisdiction even over a pending, argued case, and courts may not inquire into its motives. The repeal reached only appeals under the 1867 Act; jurisdiction under the Judiciary Act of 1789 survived.

\textbf{Reasoning.} Appellate jurisdiction comes from Article III but subject to ``such Exceptions'' as Congress makes, and this repeal was express. A repealed jurisdictional statute is treated as never having existed. \textit{Ex parte Yerger} later confirmed the 1789 route; \textit{Klein} and \textit{Boumediene} set outer limits.

\textbf{Examples.} \emph{Same-Side:} repealing direct appeals while certiorari remains. \emph{Opposite-Side:} stripping detainee habeas without a substitute. \emph{Fence-Sitter:} ending all Supreme Court review of school-prayer cases.

\textbf{Critique.} Broad readings imply plenary power over the Court's docket; Henry Hart argued exceptions cannot destroy its essential functions.

\textbf{Key Quotation.} ``Without jurisdiction, the court cannot proceed at all in any cause. Jurisdiction is power to declare the law, and, when it ceases to exist, the only function remaining to the court is that of announcing the fact and dismissing the cause.'' --- Chief Justice Chase, 74 U.S.\ at 514.

\end{document}
